\article{组合学与模 $p$ 的 Shimura 簇（据 Langlands 讲座）}{Combinatorics and Shimura Varieties mod $p$ (Based on Lectures by Langlands)}{R. E. Kottwitz}{2}{185--192}

Shimura 簇（Shimura varieties）理论中的一个主要问题，是将其 Hasse--Weil zeta 函数（Hasse--Weil zeta-functions）写成自守形式（automorphic forms）所对应的 $L$ 函数的乘积。文献 [1] 已经对某类 Shimura 簇解决了这个问题：相应的 $\Q$ 上代数群 $G$，是在映射
\[
\operatorname{Res}_{F/\Q}(B^*)\xrightarrow{\mathrm{norm}}\operatorname{Res}_{F/\Q}(F^*)
\]
下，$\operatorname{Res}_{F/\Q}(F^*)$ 的某个连通 $\Q$ 子群的逆像；这里 $F$ 是全实数域（totally real number field），$B$ 是 $F$ 上的全不定四元数代数（totally indefinite quaternion algebra）。本文只讨论 $G=\operatorname{Res}_{F/\Q}(B^*)$ 的情形。在这个特例中，相应 Shimura 簇的 zeta 函数可表示成与群 $G$ 本身相联系的自守 $L$ 函数的乘积（一般情形还需要 $G$ 以外的群）。文献 [2] 已讨论过这种情形，本文将详细说明该文末页提出的组合计算。对前述 $\operatorname{Res}_{F/\Q}(B^*)$ 的所有子群，[1] 的第 4 节已详细处理这一练习；这里在所选的特例下重做这一计算，并在适当处进一步简化假设，以便突出论证的主要思路。

先概述 [2] 如何将问题归结为这一组合计算。设 $K$ 是 $G(\A_f)$ 的紧开子群（compact open subgroup），其中 $\A_f$ 是 $\Q$ 的阿代尔环（adele ring）中无穷分量为 $0$ 的部分。设 $S_K$ 为按 [2] 的方法由 $G$ 和 $K$ 得到的 Shimura 簇。由于 $B$ 全不定，$S_K$ 定义在 $\Q$ 上，因此它的 zeta 函数 $Z(s,S_K)$ 是以 $\Q$ 的各个赋值（places）$v$ 为指标的局部因子（local factors）$Z_v(s,S_K)$ 的乘积。

设 $F'$ 是 $F$ 的有限扩域，并且在 $\Q$ 上为 Galois 扩张。对
$\sigma\in\Gal(F'/\Q)/\Gal(F'/F)$，令 $H_\sigma=\GL_2(\C)$，并令 $V_\sigma$ 为 $H_\sigma$ 在 $\C^2$ 上的标准表示（standard representation）。于是 ${}^LG^\circ=\prod_\sigma H_\sigma$ 在 $V=\bigotimes_\sigma V_\sigma$ 上有自然表示 $r$。对 $\tau\in\Gal(F'/\Q)$，规定
\[
r(\tau)(\bigotimes_\sigma v_\sigma)=\bigotimes_\sigma w_\sigma,
\qquad w_\sigma=v_{\tau^{-1}\sigma},
\]
即可把它延拓为 $L$ 群（$L$-group）${}^LG=\Gal(F'/\Q)\ltimes{}^LG^\circ$ 的表示。

[2] 的结果是，除有限多个局部因子外，
\[
Z(s,S_K)=\prod_\pi L(s-d/2,\pi,r)^{m(\pi,K)},
\]
其中 $d=[F:\Q]$，$\pi$ 遍历 $L^2(G(\Q)Z(\R)\backslash G(\A))$ 中出现的 $G(\A)/Z(\R)$ 的表示（$Z$ 为 $G$ 的中心），而 $m(\pi,K)$ 是按 [2] 所述方式与 $\pi$ 和 $K$ 联系的整数。这个关于 $\pi$ 的乘积实际上是有限的，因为除有限多个 $\pi$ 外，$m(\pi,K)$ 都是 $0$。精确的结果是
\begin{equation}
 Z_p(s,S_K)=\prod_\pi L(s-d/2,\pi_p,r)^{m(\pi,K)}.\tag{1}
\end{equation}
这里 $\pi_p$ 是 $\pi$ 在某个有限有理素数 $p$ 处的局部因子，而 $p$ 满足：
\begin{enumerate}
\item[(A)] $F$ 在 $p$ 处非分歧（unramified）；
\item[(B)] $B$ 在 $F$ 的每个位于 $p$ 之上的素位处分裂（split）；
\item[(C)] $K=K_pK^p$，其中 $K_p$ 是 $G(\Q_p)$ 的极大紧子群（maximal compact subgroup），$K^p$ 是 $G(\A_f^p)$ 的紧开子群，$\A_f^p$ 为 $\Q$ 的在 $p$ 和 $\infty$ 处分量为 $0$ 的阿代尔环。
\end{enumerate}
本文余下部分固定一个满足 (A)--(C) 的有理素数 $p$。这些假设保证 $S_K$ 在 $p$ 处有良约化（good reduction），并且除非 $\pi_p$ 非分歧，否则 $m(\pi,K)=0$。对 (1) 两边取对数，利用 $\log(1-X)$ 的幂级数展开，并比较两边 $p^{-s}$ 各次幂的系数，可见 (1) 等价于：对每个正整数 $j$，
\begin{equation}
\operatorname{Card}S_K(\F_{p^j})=\sum_\pi m(\pi,K)p^{jd/2}\operatorname{tr}(r(g_{\pi_p})^j).\tag{2}
\end{equation}
这里 $\pi$ 遍历在 $p$ 处非分歧且出现在 $L^2(G(\Q)Z(\R)\backslash G(\A))$ 中的 $G(\A)/Z(\R)$ 的全部表示，而 $g_{\pi_p}$ 是 ${}^LG$ 中与 $\pi_p$ 对应的半单共轭类（semisimple conjugacy class）的一个元素。

对 $\Q$ 的每个赋值 $v$，在 $G(\Q_v)$ 上选择 Haar 测度（Haar measure）$dg_v$，使几乎所有有限位 $v$ 处的 $G(\Z_v)$ 的测度为 $1$。用 $dg_v$ 定义 $L^1(G(\Q_v))$ 在 $G(\Q_v)$ 的表示上的作用。在 $Z(\R)$ 上选择 Haar 测度 $dz_\infty$，并用商测度 $dg_\infty/dz_\infty$ 定义 $L^1(G(\R)/Z(\R))$ 在 $G(\R)/Z(\R)$ 的表示上的作用。令 $f^p$ 为 $G(\A_f^p)$ 中子集 $K^p$ 的示性函数（characteristic function）除以 $K^p$ 关于 $\prod_{v\ne p,\infty}dg_v$ 的测度。[2] 指出，存在
$f_\infty\in C_c^\infty(G(\R)/Z(\R))$，使
\[
m(\pi,K)=m(\pi)\operatorname{tr}\pi^p(f^pf_\infty),
\]
其中 $\pi^p=\bigotimes_{v\ne p}\pi_v$，$m(\pi)$ 是 $\pi$ 在 $L^2(G(\Q)Z(\R)\backslash G(\A))$ 中的重数（multiplicity）。（已知这一重数为 $1$，但我们不必使用此事实。）由 Hecke 代数（Hecke algebras）的理论，Hecke 代数 $\mathcal H(G(\Q_p),K_p)$ 中存在唯一函数 $f_p^{(j)}$，使
\[
\operatorname{tr}\pi_p(f_p^{(j)})=
\begin{cases}
p^{jd/2}\operatorname{tr}(r(g_{\pi_p})^j),&\pi_p\text{ 非分歧},\\
0,&\pi_p\text{ 分歧}.
\end{cases}
\]
故 (2) 的右边成为 $\sum_\pi m(\pi)\operatorname{tr}\pi(f_p^{(j)}f^pf_\infty)$，借助迹公式（trace formula），可将这一谱侧求和改写为轨道积分（orbital integrals）。对 $\gamma\in G(\Q)$，设 $d\bar g$ 为 $G(\A)$ 上的 $\prod_vdg_v$ 除以 $G_\gamma(\A)$ 上某个 Haar 测度 $dg_\gamma$ 得到的商测度（如通常那样，$G_\gamma$ 表示 $\gamma$ 在 $G$ 中的中心化子）。设
\[
m_\gamma=\operatorname{meas}(G_\gamma(\Q)Z(\R)\backslash G_\gamma(\A)),
\]
这里所用的测度是 $dg_\gamma$ 除以 $G_\gamma(\Q)Z(\R)\simeq G_\gamma(\Q)\times Z(\R)$ 上的测度；后者是 $G_\gamma(\Q)$ 上赋予每点测度 $1$ 的 Haar 测度与 $dz_\infty$ 的乘积。应用迹公式，得到 (2) 等价于
\begin{equation}
\operatorname{Card}S_K(\F_{p^j})=\sum_\gamma m_\gamma
\int_{G_\gamma(\A)\backslash G(\A)}f_p^{(j)}f^pf_\infty(g^{-1}\gamma g)\,d\bar g.\tag{3}
\end{equation}
其中 $\gamma$ 遍历 $G(\Q)$ 中各共轭类的一组代表。积分
$\int_{G_\gamma(\R)\backslash G(\R)} f_\infty(g^{-1}\gamma g)\,d\bar g_\infty$
仅在 $\gamma$ 于 $\infty$ 处为中心元或椭圆元（elliptic）时才可能非零，因此 (3) 的右边为
\begin{equation}
\sum_\gamma a_\gamma\int_{G_\gamma(\A_f)\backslash G(\A_f)}
f^pf_p^{(j)}(g^{-1}\gamma g)\,d\bar g_f,\tag{4}
\end{equation}
其中
$a_\gamma=m_\gamma\int_{G_\gamma(\R)\backslash G(\R)}f_\infty(g^{-1}\gamma g)\,d\bar g_\infty$，$\gamma$ 遍历 $G(\Q)$ 中于 $\infty$ 处为椭圆元或中心元的共轭类的一组代表。

下一步需要计算 $\operatorname{Card}S_K(\F_{p^j})$。原则上，只要描述出集合 $S_K(\overline\F_p)$ 及 Frobenius 元 $\Phi$ 的作用，就能同时得到所有 $j$ 对应的点数。

\paragraph*{假设。}为简化描述，以下假设 $p$ 在 $F$ 中惰性，即 $p\mathcal O_F$ 仍为素理想。

首先，对每个三元组 $\rho=(L,m',m'')$，其中 $L$ 为 $F$ 的在 $p$ 处分裂且可嵌入 $B$ 的全虚二次扩域，$m',m''$ 为满足 $m''>m'\ge0$ 及 $m'+m''=d$ 的整数，都有 $S_K(\overline\F_p)$ 中一个相应的 $\Phi$ 稳定子集 $Y_\rho$。$S_K(\overline\F_p)$ 是这些 $Y_\rho$ 及另一个 $\Phi$ 稳定子集 $Y_0$ 的不交并。把 $L^*$ 看成 $\Q$ 上的代数群，记作 $I_\rho$。于是，作为 $\Gal(\overline\F_p/\F_p)$ 集合，$Y_\rho$ 同构于
\[
I_\rho(\Q)\backslash(G(\A_f^p)\times X_\rho)/K^p,
\]
其中 $X_\rho$ 是某个带有 $I_\rho(\Q_p)$ 和 $\Phi$ 的作用的集合，这两个作用相互交换。嵌入 $L\hookrightarrow B$ 给出映射 $I_\rho\to G$，$I_\rho(\Q)$ 经
$I_\rho(\Q)\to I_\rho(\A_f^p)\to G(\A_f^p)$ 作用于 $G(\A_f^p)$，并经 $I_\rho(\Q)\to I_\rho(\Q_p)$ 作用于 $X_\rho$；$K^p$ 则只在乘积 $G(\A_f^p)\times X_\rho$ 的第一因子上作用。

对任意正整数 $j$ 及 $x\in X_\rho$，令
$T_x^j=\{g\in I_\rho(\Q_p):\Phi^jx=gx\}$，并令 $\delta_x^j$ 为 $T_x^j$ 的示性函数。在 $I_\rho(\Q_p)$ 上选择 Haar 测度 $\mu$，对 $\gamma\in I_\rho(\Q)$ 定义
\begin{equation}
\varphi^{(j)}(\gamma)=\sum_{x\in X_\rho\bmod I_\rho(\Q_p)}\frac{1}{\mu(I_x)}\delta_x^j(\gamma),\tag{5}
\end{equation}
其中 $I_x$ 是 $x$ 在 $I_\rho(\Q_p)$ 中的稳定子。在 $G_\gamma(\A_f^p)$ 上选择 Haar 测度 $\nu$。令 $\nu'=\prod_{v\ne p,\infty}dg_v$，这里 $dg_v$ 是先前选定的 $G(\Q_v)$ 上的 Haar 测度；回忆 $f^p$ 定义为 $K^p$ 的示性函数除以 $\nu'(K^p)$。容易证明 $Y_\rho$ 对 $\operatorname{Card}S_K(\F_{p^j})$ 的贡献是
\begin{equation}
\sum_{\gamma\in I_\rho(\Q)}b_\gamma\varphi^{(j)}(\gamma)
\int_{G_\gamma(\A_f^p)\backslash G(\A_f^p)}f^p(g^{-1}\gamma g)\frac{d\nu'}{d\nu},\tag{6}
\end{equation}
其中 $b_\gamma=\operatorname{meas}(I_\rho(\Q)\backslash G_\gamma(\A_f^p)\times I_\rho(\Q_p))$。所用测度为 $\nu\times\mu$ 除以 $I_\rho(\Q)$ 上赋予每点测度 $1$ 的测度。可见 (4) 和 (6) 的形式大体相同，不过还需要对所有三元组 $\rho=(L,m',m'')$ 将 (6) 求和。(4) 中与 $m',m''$ 求和相对应的部分尚未显现。下面把 $f_p^{(j)}$ 自然地分解成若干函数之和，就能看出两者的对应。

Hecke 代数 $\mathcal H=\mathcal H(G(\Q_p),K_p)$ 中的函数 $f_p^{(j)}$，由对 $G(\Q_p)$ 的所有非分歧表示 $\pi_p$ 均成立的等式
$\operatorname{tr}\pi_p(f_p^{(j)})=p^{jd/2}\operatorname{tr}(r(g_{\pi_p})^j)$ 刻画。由于 $p$ 在 $F$ 中惰性，$F_p$ 是一个域（在 $\Q_p$ 上的次数为 $d$），并且 $G(\Q_p)\simeq\GL_2(F_p)$（因为假设了 $B$ 在 $F$ 的每个 $p$ 上方素位处分裂）。把 $\GL_2(F_p)$ 看成 $\GL_2$ 的 $F_p$ 点，可得到由 $\mathcal H$ 到关于 $a,b$ 对称的 Laurent 多项式代数 $\C[a,b,a^{-1},b^{-1}]^{\mathfrak S_2}$ 的同构 $f\mapsto f^\vee$（$a,b$ 为两个不定元）。把 $\GL_2(F_p)$ 看成 $\operatorname{Res}_{F_p/\Q_p}\GL_2$ 的 $\Q_p$ 点，则得到同构 $f\mapsto f^\sim$，其像是在 ${}^LG^\circ\times\{\Phi\}$ 上的函数代数，这些函数由复 Lie 群 ${}^LG$ 的有限维复解析表示的特征的线性组合限制而得。$f^\vee$ 与 $f^\sim$ 的关系为
\[
f^\sim\left(\begin{pmatrix}a_1&0\\0&b_1\end{pmatrix}\times\cdots\times
\begin{pmatrix}a_d&0\\0&b_d\end{pmatrix}\times\Phi\right)
=f^\vee\left(\begin{pmatrix}a_1\cdots a_d&0\\0&b_1\cdots b_d\end{pmatrix}\right).
\]
我们要找 $f_p^{(j)}\in\mathcal H$，使对所有 $x\in{}^LG^\circ\times\{\Phi\}$ 有
$(f_p^{(j)})^\sim(x)=p^{jd/2}\operatorname{tr}(r(x^j))$。令 $l$ 为 $j,d$ 的最大公因数。于是对
\[
x=\begin{pmatrix}a_1&0\\0&b_1\end{pmatrix}\times\cdots\times
\begin{pmatrix}a_d&0\\0&b_d\end{pmatrix}\times\Phi,
\]
有
\begin{align*}
\operatorname{tr}r(x^j)&=\operatorname{tr}r\left(
\begin{pmatrix}a_1\cdots a_j&0\\0&b_1\cdots b_j\end{pmatrix}
\times\begin{pmatrix}a_2\cdots a_{j+1}&0\\0&b_2\cdots b_{j+1}\end{pmatrix}
\times\cdots\times\Phi^j\right)\\
&=\big((a_1\cdots a_d)^{j/l}+(b_1\cdots b_d)^{j/l}\big)^l.
\end{align*}
因此
$(f_p^{(j)})^\vee(\diag(a,b))=p^{jd/2}(a^{j/l}+b^{j/l})^l$。
对满足 $0\le k'<l/2$ 的 $k'\in\Z$，令 $f_{j,k'}$ 为 Hecke 代数中满足
\[
f_{j,k'}^\vee=p^{jd/2}\binom l{k'}\big(a^{jk'/l}b^{jk''/l}+a^{jk''/l}b^{jk'/l}\big)
\]
的函数，其中 $k''=l-k'$。若 $l$ 为偶数，则对 $k'=k''=l/2$，令 $f_{j,k'}$ 为满足
$f_{j,k'}^\vee=p^{jd/2}\binom l{k'}a^{jk'/l}b^{jk'/l}$ 的 Hecke 代数函数。由二项式定理，
\[
f_p^{(j)}=\sum_{k'=0}^{[l/2]}f_{j,k'}.
\]

现在考虑三元组 $\rho=(L,m',m'')$，并取 $\gamma\in L-F$。回忆 $m''$ 由 $m'$ 决定（因为 $m'+m''=d$）。选择嵌入 $L\hookrightarrow B$，$\gamma$ 就给出 $I_\rho(\Q)$ 及 $G(\Q)$ 的元素。我们将证明：除非 $m'$ 可被 $d/l$ 整除，否则 $\gamma$ 对 (6) 的贡献为 $0$；可整除时，此贡献等于 $k'=lm'/d$ 和 $\gamma$ 对下面这个改写后的 (4) 的贡献：
\[
\sum_\gamma\sum_{k'=0}^{[l/2]}a_\gamma
\int_{G_\gamma(\A_f^p)\backslash G(\A_f^p)}f^p(g^{-1}\gamma g)\frac{d\nu'}{d\nu}
\int_{G_\gamma(\Q_p)\backslash G(\Q_p)}f_{j,k'}(g^{-1}\gamma g)\frac{dg_p}{dg_\gamma}.
\]
这里 $dg_\gamma$ 是 $G_\gamma(\Q_p)$ 上的 Haar 测度，在嵌入 $L\hookrightarrow B$ 诱导的同构
$I_\rho(\Q_p)\simeq G_\gamma(\Q_p)$ 下，它与 $I_\rho(\Q_p)$ 上的 Haar 测度 $\mu$ 相对应。容易证明，对我们选择的测度，$a_\gamma=b_\gamma$。所以只需证明如下结论。

\begin{originaltheorem}{定理（Theorem）}
沿用以上记号，除非 $d/l$ 整除 $m'$，否则 $\varphi^{(j)}(\gamma)=0$；可整除时，它等于
\[
\int_{G_\gamma(\Q_p)\backslash G(\Q_p)}f_{j,k'}(g^{-1}\gamma g)\frac{dg_p}{dg_\gamma},
\qquad k'=lm'/d.
\]
\end{originaltheorem}

先计算这个轨道积分。由于 $L$ 在 $p$ 处分裂，而 $f\in\mathcal H$ 的双曲轨道积分（hyperbolic orbital integrals）可直接从 $f^\vee$ 读出，这一步很容易。设 $\alpha,\beta\in F_p^*$ 为 $\gamma$ 的两个特征值，不妨假设 $\alpha$ 的赋值 $v(\alpha)$ 不大于 $\beta$ 的赋值 $v(\beta)$。则 $f_{j,k'}$ 沿 $\gamma$ 的轨道的积分为 $0$，除非 $v(\alpha)=jk'/l$ 且 $v(\beta)=jk''/l$；此时积分为
\[
\left|\frac{\alpha\beta}{(\alpha-\beta)^2}\right|_{F_p}^{1/2}
p^{jd/2}\binom l{k'}\operatorname{meas}(I_0)^{-1},
\]
其中 $I_0$ 是 $I_\rho(\Q_p)$ 的子群，在
$I_\rho(\Q_p)\simeq F_p^*\times F_p^*$ 下对应于
$\mathcal O_{F_p}^*\times\mathcal O_{F_p}^*$。最后结果是：除非
$v(\alpha)=jk'/l$ 且 $v(\beta)=jk''/l$，否则积分为 $0$；此时积分为
$\binom l{k'}p^{jk'd/l}\mu(I_0)^{-1}$。

接下来计算 $\varphi^j(\gamma)$。这就是本文开头所说的组合计算。先描述 $X_\rho$ 以及 $I_\rho(\Q_p)$ 和 $\Phi$ 在其上的作用。令 $\Q_p^{\mathrm{un}}$ 为包含 $F_p$ 的 $\Q_p$ 的极大非分歧扩张，$\mathcal O_p^{\mathrm{un}}$ 为其赋值环（valuation ring），$\mathcal B'$ 为 $\Q_p^{\mathrm{un}}$ 上二维向量空间
$\Q_p^{\mathrm{un}}\oplus\Q_p^{\mathrm{un}}$ 中所有 $\mathcal O_p^{\mathrm{un}}$ 格（lattices）的集合，$\mathcal B$ 为其中格的等价类的集合（若存在非零 $c\in\Q_p^{\mathrm{un}}$ 使 $M_1=cM_2$，就称 $M_1,M_2$ 同类）。对任意格 $M$，用 $\overline M$ 表示其类。$\mathcal B$ 是 $\SL_2(\Q_p^{\mathrm{un}})$ 的 Bruhat--Tits 建筑（Bruhat--Tits building）的顶点集合。这座建筑是一棵树。群 $\GL_2(\Q_p^{\mathrm{un}})$ 和 $\Gal(\Q_p^{\mathrm{un}}/\Q_p)$ 均作用于 $\mathcal B'$ 和 $\mathcal B$。令
\[
B=\begin{pmatrix}p^{m'}&0\\0&p^{m''}\end{pmatrix},
\]
并令 $\sigma$ 为 $\Gal(\Q_p^{\mathrm{un}}/\Q_p)$ 中的 Frobenius 元。令 $X_\rho$ 为由 $\mathcal B'$ 的元素 $M_i$ 组成的双无限序列 $\{M_i\}_{i\in\Z}$ 的集合，要求
\begin{enumerate}
\item[(A)] 对每个 $i\in\Z$，$M_i\supsetneqq M_{i-1}\supsetneqq pM_i$；
\item[(B)] 对每个 $i\in\Z$，$B\sigma^dM_i=M_{i-d}$。
\end{enumerate}
$\Phi$ 在 $X_\rho$ 上的作用为 $\Phi:\{M_i\}\mapsto\{M_i'\}$，其中 $M_i'=M_{i-1}$。把 $I_\rho(\Q_p)$ 与 $A(F_p)$ 等同，其中 $A$ 是对角环面。元素 $\gamma\in I_\rho(\Q_p)$ 在 $X_\rho$ 上的作用由 $\gamma\{M_i\}=\{\gamma M_i\}$ 给出。

回忆我们要计算
$\varphi^j(\gamma)=\sum_x\mu(I_x)^{-1}\delta_x^j(\gamma)$，其中对 $x\in X_\rho\bmod I_\rho(\Q_p)$ 求和。令 $I_1$ 为 $I_\rho(\Q_p)=A(F_p)$ 中所有形如
\[
\begin{pmatrix}cp^n&0\\0&c\end{pmatrix},\qquad n\in\Z,\quad c\in F_p^*,
\]
的矩阵组成的子群。仍令 $I_0=A(\mathcal O_{F_p})$，则
$I_\rho(\Q_p)=I_0\times I_1$。\footnote{\textbf{译者注（本次校订）：}按这里对 $I_1$ 的定义，此处应为子群的乘积 $I_\rho(\Q_p)=I_0I_1$，而非直积。两子群的交是标量单位矩阵群 $\mathcal O_{F_p}^*1_2$。这些标量固定每个格，因而属于所有 $I_x$；所以后面的指数等式 $[I_0:I_x]=[I_\rho(\Q_p):I_1I_x]$ 仍然成立。}此外，
\[
\mu(I_0)=\mu(I_x)[I_0:I_x]=\mu(I_x)[I_\rho(\Q_p):I_1I_x],
\]
故 $\varphi^j(\gamma)=\mu(I_0)^{-1}\operatorname{Card}S(\gamma,j)$，其中
\[
S(\gamma,j)=\{x\in X_\rho\bmod I_1:\Phi^jx=\gamma x\}.
\]
设 $\mathfrak A$ 为 $\mathcal B$ 中对应于分裂环面 $A$ 的公寓（apartment）。在 $\mathfrak A$ 中取一点 $p_0$，并令
\[
\mathcal M=\{x\in\mathcal B:p_0\text{ 是 }\mathfrak A\text{ 中离 }x\text{ 最近的点}\}.
\]
令 $\mathcal M'$ 为 $\mathcal M$ 中各格类的一组代表；换言之，$\mathcal M'$ 是 $\mathcal B'$ 的一个子集，其中没有两个元素同类，且
$\{\overline M:M\in\mathcal M'\}=\mathcal M$。于是 $X_\rho$ 中满足
\begin{enumerate}\item[(C)] $M_0\in\mathcal M'$\end{enumerate}
的全部 $\{M_i\}$ 构成 $X_\rho\bmod I_1$ 的一组代表。条件
$\Phi^j\{M_i\}=\gamma\{M_i\}$ 恰好就是
\begin{enumerate}\item[(D)] $\gamma M_i=M_{i-j}$。\end{enumerate}
综上，$S(\gamma,j)$ 是 $\mathcal B'$ 的元素组成的、满足 (A)--(D) 的所有序列的集合。$S(\gamma,j)$ 中每个序列 $\{M_i\}$ 决定 $\mathcal B$ 中的序列 $\{x_i\}=\{\overline M_i\}$，它满足：
\begin{enumerate}
\item[(A$'$)] $x_i$ 是 $x_{i-1}$ 的邻点；
\item[(B$'$)] $B\sigma^dx_i=x_{i-d}$；
\item[(C$'$)] $x_0\in\mathcal M$；
\item[(D$'$)] $\gamma x_i=x_{i-j}$。
\end{enumerate}
令 $S'(\gamma,j)$ 为满足 (A$'$)--(D$'$) 的序列 $\{x_i\}$ 的集合。若 $\det\gamma$ 的赋值不等于 $j$，则 (A) 和 (D) 不相容，$S(\gamma,j)$ 为空。因此从现在起假设 $v(\det\gamma)=j$。在此情形下，$S'(\gamma,j)$ 的任意元素都可唯一地提升为 $S(\gamma,j)$ 的元素（为说明这一点，须使用 $v(\det B)=d$）。故当 $v(\det\gamma)=j$ 时，
$\operatorname{Card}S(\gamma,j)=\operatorname{Card}S'(\gamma,j)$。取 $r,s\in\Z$ 使 $rj+sd=l$。条件 (B$'$) 和 (D$'$) 合起来等价于：
\begin{enumerate}
\item[(E$'$)] $\sigma^{dj/l}x_i=B^{-j/l}\gamma^{d/l}x_i$；
\item[(F$'$)] $\gamma^rB^s\sigma^{sd}x_i=x_{i-l}$。
\end{enumerate}
现在证明：除非 $B^{-j/l}\gamma^{d/l}\in I_0$，否则 $S'(\gamma,j)$ 为空。设其非空，取 $\{x_i\}\in S'(\gamma,j)$。因为 $p_0$ 是 $\mathfrak A$ 中离 $x_0$ 最近的点，$\tau p_0$ 就是 $\tau\mathfrak A$ 中离 $\tau x_0$ 最近的点，其中
$\tau=\sigma^{-dj/l}B^{-j/l}\gamma^{d/l}$。由 (E$'$)，$\tau x_0=x_0$；又因为
$\sigma\mathfrak A=\mathfrak A$、$B\mathfrak A=\mathfrak A$ 及 $\gamma\mathfrak A=\mathfrak A$，有 $\tau\mathfrak A=\mathfrak A$。因此 $\tau$ 固定 $p_0$。$\sigma^{-dj/l}$ 也固定 $p_0$，故 $B^{-j/l}\gamma^{d/l}$ 固定 $p_0$。但是这个矩阵是对角矩阵，且行列式为单位（因为假设了 $v(\det\gamma)=j$），所以固定 $p_0$ 就意味着它属于 $I_0$。这就证明了相应的消失结论。

设 $\alpha,\beta\in F_p^*$ 为 $\gamma$ 的对角元，故 $\gamma=\diag(\alpha,\beta)$。那么
\[
B^{-j/l}\gamma^{d/l}\in I_0
\quad\Longleftrightarrow\quad
v(\alpha)=jm'/d,\quad v(\beta)=jm''/d.
\]
特别地，除非 $jm'/d,jm''/d$ 均为整数，否则对所有 $\gamma$ 都有 $\varphi^j(\gamma)=0$。回忆 $l=(j,d)$ 及 $m'+m''=d$，可见 $jm'/d,jm''/d\in\Z$ 当且仅当 $d/l$ 整除 $m'$；即使可整除，除非
\[
v(\alpha)=jm'/d=jk'/l,\qquad v(\beta)=jm''/d=jk''/l,
\]
$\varphi^j(\gamma)$ 仍为零。与此前算出的轨道积分比较，可见只剩下如下引理需要证明。

\begin{originaltheorem}{引理（Lemma）}
设 $B^{-j/l}\gamma^{d/l}\in I_0$。则
$\operatorname{Card}S'(\gamma,j)=\binom l{k'}p^{jk'd/l}$，其中 $k'=lm'/d$。
\end{originaltheorem}

为证明此引理，先仔细考察 (E$'$) 和 (F$'$)。令
$\mathcal B^\tau=\{x\in\mathcal B:\tau(x)=x\}$，其中 $\tau$ 是此前定义的 $\mathcal B$ 的自同构。于是 (E$'$) 就是说对所有 $i\in\Z$，$x_i\in\mathcal B^\tau$。令 $\Q_p^c$ 为 $\Q_p^{\mathrm{un}}$ 的完备化。存在 $\GL_2(\Q_p^c)$ 在 $\mathcal B$ 上的唯一连续作用，延拓 $\GL_2(\Q_p^{\mathrm{un}})$ 在 $\mathcal B$ 上的作用（$\mathcal B$ 赋予离散拓扑）。$\GL_2(\Q_p^c)$ 的每个元素都以树 $\mathcal B$ 的单纯自同构（simplicial automorphism）作用。利用 $B^{-j/l}\gamma^{d/l}\in I_0$，容易证明存在对角元属于 $\Q_p^c-\{0\}$ 的对角矩阵 $\delta$，使
\[
\delta^{\sigma^{dj/l}}\delta^{-1}=B^{-j/l}\gamma^{d/l}.
\]
简单计算表明
$\mathcal B^\tau=\delta\mathcal B^{\sigma^{dj/l}}$，其中
$\mathcal B^{\sigma^{dj/l}}$ 表示 $\sigma^{dj/l}$ 在 $\mathcal B$ 中的不动点集合。这个子树可典范地等同于 $\Q_p^{\mathrm{un}}$ 在 $\sigma^{dj/l}$ 下的不动域上的 $\SL_2$ 的 Bruhat--Tits 建筑。因此它是一棵每个顶点恰有 $p^{dj/l}+1$ 个邻点的树。由于 $\mathcal B^\tau$ 只是它经 $\delta$ 的平移，它也有同样的性质。

$B^{-j/l}\gamma^{d/l}\in I_0$ 还推出 $\mathcal B^\tau$ 包含 $\mathfrak A$；这是因为 $\sigma^{dj/l}$ 以及 $I_0$ 的每个元素均逐点固定 $\mathfrak A$。

(F$'$) 中出现了 $\mathcal B$ 的自同构 $\gamma^rB^s\sigma^{sd}$。我们需要知道它在 $\mathfrak A$ 上的作用。$\sigma^{sd}$ 在 $\mathfrak A$ 上平凡作用；$B$ 在 $\mathfrak A$ 上作距离为 $m''-m'$ 的平移。因为 $B^{-j/l}\gamma^{d/l}\in I_0$，$\gamma^d$ 在 $\mathfrak A$ 上的作用与 $B^j$ 相同，所以 $\gamma$ 作距离为 $(m''-m')j/d$ 的平移。因此 $\gamma^rB^s\sigma^{sd}$ 在 $\mathfrak A$ 上作距离为 $l-2k'$ 的平移。这样，剩下的工作就是证明下面的树上计数引理。

\begin{originaltheorem}{引理（Lemma）}
设 $T$ 是每个顶点有 $q+1$ 个邻点的树，$\mathfrak A$ 是 $T$ 的一个子树，其中每个顶点有 $2$ 个邻点。取 $p_0\in\mathfrak A$。设 $l,k$ 是满足 $l-2k>0$ 的非负整数，$\xi$ 是 $T$ 的自同构，在 $\mathfrak A$ 上作距离为 $l-2k$ 的平移。令 $N(l,k)$ 表示由 $T$ 中 $l+1$ 个点 $x_0,\ldots,x_l$ 组成的、满足下列条件的序列的数目：
\begin{enumerate}
\item[(i)] 对 $i=0,\ldots,l-1$，$x_i$ 是 $x_{i+1}$ 的邻点；
\item[(ii)] $p_0$ 是 $\mathfrak A$ 中离 $x_0$ 最近的点；
\item[(iii)] $x_l=\xi x_0$。
\end{enumerate}
则 $N(l,k)=\binom lkq^k$。
\end{originaltheorem}

约定把满足 (i) 的点列 $x_0,\ldots,x_l$ 称为连接 $x_0$ 与 $x_l$ 的长 $l$ 的路径（path）。于是
\[
N(l,k)=\sum_{x_0\in\mathcal M}\operatorname{Card}\{\text{连接 }x_0\text{ 与 }\xi x_0\text{ 的长 }l\text{ 的路径}\},
\]
其中 $\mathcal M$ 为满足 (ii) 的点 $x_0\in T$ 的集合。

设 $a,b\in\Z$ 且 $a\ge2b\ge0$，$x,y$ 是 $T$ 中距离为 $a-2b$ 的两个顶点。那么连接 $x,y$ 的长 $a$ 的路径数与 $x,y$ 的选择无关，记为 $M(a,b)$。

下图表明，$x_0$ 与 $\xi x_0$ 之间的距离等于 $l-2k+2n$，其中 $n$ 为 $x_0$ 到 $p_0$ 的距离。
\begin{center}
\begin{tikzpicture}[scale=.85]
\draw (-1.2,0)--(7,0);\node[right] at (7,0){$\mathfrak A$};
\draw (0,0)--(0,2.1);\draw (5,0)--(5,2.1);
\foreach \x in {-1,0,1,2,3,4,5,6}{\fill (\x,0) circle (1.5pt);}
\foreach \x in {0,5}{\fill (\x,1.05) circle (1.5pt);\fill (\x,2.1) circle (1.5pt);}
\node[below] at (0,0){$p_0$};\node[below] at (5,0){$\xi p_0$};
\node[above right] at (0,2.1){$x_0$};\node[above right] at (5,2.1){$\xi x_0$};
\draw[<->] (-.4,.1)--node[left]{$n$}(-.4,2);
\draw[<->] (5.45,.1)--node[right]{$n$}(5.45,2);
\draw[<->] (.2,-.65)--node[below]{$l-2k$}(4.8,-.65);
\end{tikzpicture}
\end{center}
因此
\[
N(l,k)=\sum_{n=0}^k\operatorname{Card}\{x_0\in\mathcal M:\text{$x_0$ 到 $p_0$ 的距离为 $n$}\}\,M(l,k-n).
\]
到 $p_0$ 距离为 $n$，且以 $p_0$ 为 $\mathfrak A$ 中最近点的点 $x_0$，在 $n>0$ 时有 $(1-q^{-1})q^n$ 个，在 $n=0$ 时有 $1$ 个。因此
\[
N(l,k)=M(l,k)+(1-q^{-1})\sum_{n=1}^kq^nM(l,k-n).
\]
函数 $M(l,k)$ 定义在满足 $l\ge2k\ge0$ 的整数对 $(l,k)$ 上，由以下三个性质刻画：
\begin{enumerate}
\item[(i)] 对 $l\ge0$，$M(l,0)=1$；
\item[(ii)] 当 $l=2k>0$ 时，$M(l,k)=(q+1)M(l-1,k-1)$；
\item[(iii)] 当 $l>2k>0$ 时，$M(l,k)=M(l-1,k)+qM(l-1,k-1)$。
\end{enumerate}
(i) 直接成立，(i)--(iii) 也显然能递归地确定 $M$。为验证 (ii)、(iii)，只需按路径的第一步分类：先选择 $x$ 的邻点 $x'$，再数从 $x'$ 到 $y$ 的长 $l-1$ 的路径。

对 $l,k\in\Z$，令 $m(l,k)$ 为
\[
f(X,Y)=(1-qY)(1-Y)^{-1}(1-X-qXY)^{-1}
\]
的形式幂级数（formal power series）展开中 $X^lY^k$ 的系数。我们将证明：当 $l\ge2k\ge0$ 时，$M(l,k)=m(l,k)$。只需说明将 $M$ 换成 $m$ 后 (i)--(iii) 仍成立。由 $f(X,Y)$ 的定义，
$(1-X-qXY)f(X,Y)=(1-qY)(1-Y)^{-1}$。比较两边 $X^lY^k$ 的系数，得到当 $l>0$ 时，
\[
m(l,k)-m(l-1,k)-qm(l-1,k-1)=0,
\]
这说明 $m$ 满足 (iii)。在 $f$ 的定义中令 $Y=0$，得到
$\sum_{l=0}^\infty m(l,0)X^l=(1-X)^{-1}=1+X+X^2+\cdots$，说明 $m$ 满足 (i)。

还需证明 $m$ 满足 (ii)。令 $g(X,Y)=(1-X-qXY)^{-1}$，并令 $n(l,k)$ 为其形式幂级数展开中 $X^lY^k$ 的系数。有
\begin{align*}
(1-X-qXY)^{-1}
&=\sum_{l=0}^\infty(X+qXY)^l
=\sum_{l=0}^\infty\sum_{k=0}^\infty\binom lkX^{l-k}(qXY)^k\\
&=\sum_{l\ge0,\,k\ge0}\binom lkq^kX^lY^k,
\end{align*}
故 $n(l,k)=\binom lkq^k$。特别地，对 $k>0$，
\[
n(2k-1,k)=\binom{2k-1}{k}q^k
=\binom{2k-1}{k-1}q^{k-1}q=qn(2k-1,k-1).
\]
这表明 $(1-qY)g(X,Y)$ 中 $X^{2k-1}Y^k$ 的系数为零。但
$(1-qY)g(X,Y)=(1-Y)f(X,Y)$，故后者中这个系数也为零，也就是说
\begin{equation}
m(2k-1,k)=m(2k-1,k-1)\quad\text{对一切 }k>0.\tag{7}
\end{equation}
我们已经知道，当 $l>0$ 时，$m(l,k)=m(l-1,k)+qm(l-1,k-1)$。取 $l=2k$，$k>0$，则
$m(2k,k)=m(2k-1,k)+qm(2k-1,k-1)$。与 (7) 合并，即知 $m$ 满足 (ii)。

回到 $N(l,k)$ 的讨论，可见
\[
N(l,k)=m(l,k)+(1-q^{-1})\sum_{n=1}^\infty q^nm(l,k-n),
\]
因为只要 $M$ 有定义，$m$ 就与 $M$ 一致，而当 $n>k$ 时 $m(l,k-n)=0$。可用这个等式把 $N(l,k)$ 的定义延拓到所有整数 $l,k$。按此约定，$\sum N(l,k)X^lY^k$ 等于
\begin{align*}
&\sum_{l,k}m(l,k)X^lY^k
 +(1-q^{-1})\sum_{n=1}^\infty q^n\sum_{l,k}m(l,k-n)X^lY^k\\
&=f(X,Y)+(1-q^{-1})\sum_{n=1}^\infty q^nY^n
 \sum_{l,k}m(l,k-n)X^lY^{k-n}\\
&=f(X,Y)\left[1+(1-q^{-1})\sum_{n=1}^\infty q^nY^n\right]\\
&=f(X,Y)(1-Y)(1-qY)^{-1}=g(X,Y).
\end{align*}
$g(X,Y)$ 的形式幂级数展开已在前面算出。比较系数即得 $N(l,k)=\binom lkq^k$，引理得证。

\begin{thebibliography}{2}
\bibitem[1]{c31-1} R. P. Langlands, \textit{On the zeta-functions of some simple Shimura varieties}, Notes, Institute for Advanced Study, Princeton, N. J., 1977.
\bibitem[2]{c31-2} R. P. Langlands, \textit{Shimura varieties and the Selberg trace formula}, Canad. J. Math.（待刊）。
\end{thebibliography}
\noindent 美国华盛顿大学，西雅图（University of Washington, Seattle）。
\footnotetext{AMS（MOS）主题分类（1970）：05C05、14G10、22E50、22E55。原文版权：1979 年，美国数学会。}
